\documentclass{article}

\usepackage[preprint]{tmlr}

\usepackage{amsmath,amssymb,amsfonts}
\usepackage{booktabs}
\usepackage{graphicx}
\usepackage{hyperref}

\title{Neural Geometry Is Not Metric-Neutral: \\
Dimensionality Mediates Dissociation Between Similarity Metrics in Brain-Wide Neuropixels Recordings}

\author{\name Elliot Tower \email elliot@elliottower.ai \\
      \addr Independent Researcher}

\begin{document}
\maketitle

\begin{abstract}
We study how the choice of geometric similarity metric interacts with brain region dimensionality. Using Neuropixels recordings across 73 brain regions in 10 mice \citep{steinmetz2019distributed}, we show that linear (CKA) and nonlinear (UMAP Procrustes) similarity are strongly anti-correlated ($\rho = -0.85$, 95\% CI $[-0.87, -0.83]$), and that this dissociation is fully mediated by effective dimensionality: controlling for the power-law exponent $\alpha$ reverses the sign to $r = +0.44$. Singular value spectra are highly conserved across animals (mean $r = 0.94$), yet temporal modes are idiosyncratic, indicating universal dimensionality \emph{structure} but animal-specific \emph{content}. Dimensionality is not fixed: under hard task conditions, $\alpha$ increases while manifold shape is preserved, revealing dimensionality compression as an active computational strategy. Evidence and choice subspaces show two encoding strategies mediated by dimensionality: low-dimensional regions route through shared subspaces, while high-dimensional regions maintain separate, flexibly rotatable representations. A potent/null space decomposition confirms the mechanism: the geometric dissociation index is strongly anti-correlated with null-space fraction ($\rho = -0.86$, $p = 5 \times 10^{-22}$), meaning regions where metrics disagree most are precisely those with the largest choice-potent subspaces. The CKA--Procrustes dissociation replicates on the independent IBL Brain-Wide Map dataset ($\rho = -0.94$, 139 mice, 12 labs). These findings establish that the choice of similarity metric is not neutral---it interacts systematically with dimensionality to determine which brain regions appear functionally similar. A companion paper addresses causal subspace discovery and the validation of nonlinear methods against optogenetic ground truth.
\end{abstract}

\section{Introduction}

A growing body of work characterizes neural computation through the geometry of population activity \citep{cunningham2014dimensionality, gallego2017neural}. When $N$ neurons fire together, their joint activity traces trajectories in an $N$-dimensional space, and the low-dimensional structure of those trajectories reveals the underlying computation. Decision variables live in low-dimensional subspaces \citep{steinmetz2019distributed}. Motor plans occupy orthogonal manifolds \citep{gallego2017neural}. Representations drift over time within fixed-dimensional structures \citep{driscoll2022representational}.

This geometric view raises a methodological question with substantive consequences: \emph{does the choice of similarity metric change the scientific conclusion?} The field's standard tools for comparing neural representations---centered kernel alignment (CKA), representational similarity analysis (RSA), and their variants---operate on kernel or covariance matrices and are therefore sensitive to the full variance structure of the population. Nonlinear methods such as UMAP Procrustes instead compare the shape of low-dimensional manifold embeddings. If these two classes of methods disagree, and if the disagreement is systematic rather than noisy, then the choice of metric is not neutral: it determines which brain regions appear functionally similar.

We test this directly. Using 1,316 cross-session pairs across 73 brain regions in the Steinmetz et al.\ (2019) Neuropixels dataset, we establish a geometric dissociation between metric families:

\textbf{Metrics disagree, and dimensionality explains why.} Linear (CKA) and nonlinear (UMAP Procrustes) similarity are strongly anti-correlated ($\rho = -0.85$). This is mediated by effective dimensionality: high-dimensional regions appear dissimilar to linear methods but highly similar to nonlinear ones. Once dimensionality is controlled for, the metrics agree (partial $r = +0.44$). This empirically validates the theoretical prediction of \citet{harvey2024representational} that CKA sensitivity depends on participation ratio.

A companion paper establishes a second dissociation: linear subspace methods systematically underestimate causal structure, validated against optogenetic ground truth.

These findings converge on a single message: \emph{the choice of geometric similarity metric interacts with brain region dimensionality}. We interpret this through the neural manifold framework \citep{gallego2017neural}: the choice subspace is the \emph{potent space} for decision output \citep{kaufman2014cortical}, and high-dimensional regions harbor large null spaces that inflate CKA dissimilarity while preserving manifold-level structure. A potent/null space decomposition confirms this: the geometric dissociation index is strongly anti-correlated with null-space fraction ($\rho = -0.86$), meaning regions where metrics disagree most are those with the largest choice-potent subspaces. The dissociation replicates on the independent IBL Brain-Wide Map dataset ($\rho = -0.94$, 139 mice, 12 labs).

\section{Related Work}

\textbf{Neural manifolds and population-level computation.} \citet{gallego2017neural} argued that the fundamental unit of motor cortical computation is not the single neuron but the \emph{neural mode}---a population-level activation pattern spanning a low-dimensional manifold. Individual corticomotoneuronal cells need not encode any specific movement covariate; what matters is the trajectory through the space of modes. \citet{kaufman2014cortical} formalized this as the potent/null space distinction: preparatory activity lives in a null space orthogonal to the muscle-driving output dimensions, while execution activity enters the potent space. \citet{sadtler2014neural} demonstrated that subjects easily adapt to perturbations within the existing neural manifold but struggle with perturbations that require new modes, confirming that the manifold structure constrains behavior. Our evidence-choice subspace analysis directly tests this framework: the choice subspace is the potent space for decision output, and high-dimensional regions may harbor large null spaces that contain variance but no causal signal.

\textbf{Representational similarity.} CKA \citep{kornblith2019similarity} and RSA \citep{kriegeskorte2008representational} are the standard tools for comparing neural representations. CKA operates on kernel matrices and is therefore linear in the population activity; RSA abstracts from neuron identity by comparing representational dissimilarity matrices (RDMs), preserving relational structure across populations of different size \citep{kriegeskorte2008representational}. \citet{williams2021generalized} showed that many similarity measures are special cases of a generalized shape metric, but all within the linear/kernel family. Our Grassmannian analysis is complementary: RSA compares the global representational structure across all stimuli, while Grassmannian distance compares the subspace structure of the choice-relevant component specifically.

\textbf{Nonlinear methods.} UMAP \citep{mcinnes2018umap} and t-SNE \citep{vandermaaten2008visualizing} are widely used for visualization but less commonly for quantitative comparison. Procrustes analysis on UMAP embeddings has been used for cross-species comparison \citep{deitch2021representational}.

\textbf{Topological and dynamical methods.} Persistent homology has been applied to detect ring-like structures in head direction cells \citep{gardner2022toroidal} and grid cells \citep{rybakken2019decoding}. jPCA \citep{churchland2012neural} extracts rotational dynamics; DSA \citep{ostrow2023beyond} compares dynamical systems.

\textbf{Dimensionality and the cortical hierarchy.} Effective dimensionality scales unboundedly with neuron number \citep{stringer2024unbounded} and changes systematically along the cortical hierarchy \citep{huntenburg2024categorical}. \citet{harvey2024representational} proved that CKA measures average alignment between optimal linear readouts, with sensitivity depending on participation ratio---predicting our empirical finding that high-dimensional regions show low CKA even when manifold shape is conserved.

\textbf{Multimodal parcellation and convergent validity.} \citet{glasser2016multi} defined 180 cortical areas per hemisphere using converging evidence from architecture, function, connectivity, and topography. This multimodal convergence is the neuroscience gold standard for validating that region boundaries are biologically real. Our geometric clustering partially recovers anatomical organization (ARI $= 0.06$, NMI $= 0.30$), suggesting that our metrics capture functional rather than architectonic boundaries.

\textbf{Task transfer structure.} \citet{zamir2018taskonomy} built a computational taxonomy of visual tasks, showing that perceptual and geometric tasks cluster by computational type and that transfer structure is asymmetric. Our cross-region Grassmannian distance matrix is a neural analogue: regions whose choice subspaces are close on the Grassmannian have learned geometrically compatible representations, and the structure of this distance matrix may reveal a functional hierarchy from sensory to association to motor regions.

\textbf{What is missing.} No prior work systematically compares metric families on the same neural data, characterizes how the choice of metric interacts with brain region dimensionality, or decomposes the potent/null space variance structure across dozens of brain regions simultaneously.

\section{Methods}

\subsection{Dataset}

We use the Steinmetz et al.\ (2019) dataset: Neuropixels recordings from 10 mice performing a two-alternative visual decision task. The dataset covers 73 brain regions with $\geq 15$ neurons per region (50 with $\geq 2$ recording sessions). We extract trial-averaged activity in a post-stimulus window (150--350ms) and compute binary choice labels (left vs.\ right). All cross-session pairs for the same region constitute the comparison set: 1,316 pairs across 50 multi-session regions.

\subsection{Geometric Invariants}

\textbf{Linear similarity (CKA).} Centered kernel alignment \citep{kornblith2019similarity} on trial $\times$ neuron activity matrices, using the linear kernel.

\textbf{Nonlinear manifold similarity (UMAP Procrustes).} 5-dimensional UMAP embedding \citep{mcinnes2018umap} with Procrustes alignment between sessions.

\textbf{Subspace distance (Grassmannian).} $k$-dimensional LDA subspaces with geodesic distance on $\operatorname{Gr}(k, n)$.

\textbf{Topological invariants.} Vietoris-Rips persistence diagrams via Ripser \citep{bauer2021ripser}.

\textbf{Dynamical invariants.} Simplified jPCA on choice-difference trajectories.

\textbf{Power-law exponent ($\alpha$).} Fit to PCA eigenvalue spectrum in log-log space (ranks 10--50), characterizing dimensionality structure.

\section{Results}

\subsection{Dissociation 1: Linear and Nonlinear Metrics Anti-Correlate}

Across 1,316 cross-session pairs and 50 brain regions, CKA and UMAP Procrustes are strongly anti-correlated: Spearman $\rho = -0.85$, 95\% CI $[-0.87, -0.83]$. Motor cortex (MOs) exemplifies the extreme: CKA $= 0.07$ (near-zero linear similarity) but UMAP Procrustes $= 0.98$ (near-maximal nonlinear similarity).

\textbf{The dissociation is robust.} Bootstrap 95\% CI: $[-0.867, -0.830]$. Leave-one-mouse-out jackknife: $\rho \in [-0.872, -0.835]$. Leave-one-region-out: $\rho \in [-0.863, -0.841]$. Neuron sub-sampling: $\rho = -0.81$ ($p < 10^{-7}$). Spectral denoising (Marchenko--Pastur thresholding) \emph{strengthens} the result: $\rho = -0.889$.

\textbf{Dimensionality is the mechanism.} Partial correlation controlling for power-law exponent $\alpha$ reverses the sign: partial $r = +0.44$. Controlling for neuron count alone gives partial $r = +0.13$. Once dimensionality is accounted for, CKA and Procrustes \emph{agree}. This empirically validates the theoretical prediction of \citet{harvey2024representational}.

\textbf{The anti-correlation is metric-specific.} Sliced Wasserstein distance is uncorrelated with dimensionality ($\rho = 0.02$, $p = 0.47$), while SPD Riemannian distance tracks it ($\rho = 0.44$). The dissociation is not a generic ``all metrics disagree'' phenomenon but a specific interaction between kernel-based and manifold-based methods mediated by spectral structure.

\subsection{Spectral Universality and Dynamic Dimensionality}

Singular value spectra are highly conserved across animals (mean $r = 0.94 \pm 0.06$), but temporal modes are idiosyncratic (mean cosine $= 0.09$). The dimensionality \emph{structure} is universal; the \emph{content} is animal-specific.

Dimensionality is not fixed. Under hard task conditions (low contrast), $\alpha$ increases by $+12.4$ (Wilcoxon $p = 3.0 \times 10^{-7}$) while manifold shape is preserved (Procrustes $= 0.90$). Dimensionality compression is an active computational strategy, not a fixed anatomical property.

Rotational dynamics appear in 72/73 regions above a temporal-shuffle null, but strength varies 50$\times$ (range $0.008$--$0.40$). Persistent homology reveals almost no topological structure: only 9/73 regions have $\beta_1 > 0$.

\subsection{Evidence-Choice Subspace Alignment}

For each region, we estimate evidence and choice subspaces and measure their Grassmannian distance. Two strategies emerge:

\begin{enumerate}
\item \textbf{Low-dimensional regions} maintain stable evidence-choice alignment across conditions (high overlap, small choice subspace shift under evidence change).
\item \textbf{High-dimensional regions} show large choice subspace rotations when evidence is experimentally changed ($\rho = -0.59$, permutation $p < 10^{-6}$), maintaining separate, flexibly rotatable representations.
\end{enumerate}

Cross-condition (easy vs.\ hard) subspace rotation also anti-correlates with $\alpha$ ($\rho = -0.65$, 95\% CI $[-0.76, -0.49]$, permutation $p < 0.002$).

\subsection{Cross-Dataset Replication: IBL Brain-Wide Map}
\label{sec:ibl}

To test whether the CKA--Procrustes dissociation generalizes beyond a single dataset, we replicated the analysis on the IBL Brain-Wide Map \citep{ibl2025brainwide}, an independent multi-laboratory dataset comprising 139 mice across 12 institutions. After matching brain regions between datasets (requiring $\geq 3$ sessions and $\geq 15$ neurons per region), 11 regions were available for direct comparison: CA1, CA3, DG, GPe, LP, POST, RSPd, SUB, VISam, VISl, and VISp.

The CKA--Procrustes anti-correlation replicates strongly in the IBL data ($\rho = -0.94$, $p = 5.5 \times 10^{-5}$; permutation test $p < 10^{-4}$, $n = 10{,}000$ permutations). The effect size is comparable to the original Steinmetz result ($\rho = -0.90$ in the matched subset, $-0.85$ across all 50 regions), despite substantial differences in recording pipeline, spike sorting, and session protocols across the 12 IBL labs. Error bars ($\pm 1$ SE across sessions per region) are modest, confirming that within-region variability does not drive the correlation.

Geometric type classification transfers for 7 of 11 matched regions (64\%). The four mismatches (CA3, DG, RSPd, VISl) are all hippocampal or higher visual regions near the CKA--Procrustes boundary, where small shifts in absolute metric values can flip the binary classification while preserving the overall anti-correlation.

Absolute power-law exponents do \emph{not} transfer across datasets ($\rho = -0.27$, n.s.), indicating that the \emph{relative} geometric relationship between metrics is robust while absolute dimensionality estimates are sensitive to recording pipeline differences (neuron yield, spike sorting, session count). This dissociation between transferable relationships and non-transferable absolutes is itself informative: it suggests that the CKA--Procrustes anti-correlation reflects a genuine property of neural population geometry rather than an artifact of any particular recording setup.

\section{Discussion}

\subsection{The Metric Disagreement Is Not Noise}

The CKA--Procrustes anti-correlation ($\rho = -0.85$) is not a nuisance or a measurement artifact. It reflects a genuine interaction between metric family and brain region dimensionality. Kernel-based methods (CKA) are dominated by the full covariance structure, which in high-dimensional regions is spread across many dimensions; manifold-based methods (UMAP Procrustes) instead capture the shape of the low-dimensional embedding, which is conserved across sessions even when the high-dimensional variance structure differs.

This creates a systematic blind spot for any study that relies on a single metric family. A researcher using CKA alone would conclude that motor and hippocampal regions have low cross-session similarity, while a researcher using Procrustes alone would conclude the opposite. Neither conclusion is wrong per se---they are answering different geometric questions---but conflating them leads to contradictory inferences about which regions maintain stable representations. The dimensionality diagnostic ($\alpha$) resolves this: once the spectral structure is known, the expected direction of disagreement is predictable. A companion paper extends this finding to show that the dimensionality interaction also affects causal subspace discovery: linear methods for identifying task-relevant subspaces (LDA) fail most severely in precisely the high-dimensional regions where the metric dissociation is largest.

\subsection{The Neural Manifold Hypothesis as Unifying Framework}

Our results are naturally interpreted through the neural manifold framework of \citet{gallego2017neural}. The choice-relevant subspace corresponds to the \emph{potent space} for decision output \citep{kaufman2014cortical}---the subspace of population activity that drives the downstream binary decision. The remaining variance, which may be high-dimensional in association regions like MOs, lives in the \emph{null space}---it is present, it contributes to CKA dissimilarity, but it does not drive behavior.

This potent/null decomposition explains the CKA--Procrustes dissociation. A region like MOs shows CKA $\approx 0$ because the high-dimensional null space dominates the kernel matrix, making sessions appear dissimilar. But UMAP Procrustes $\approx 1$ because the manifold \emph{shape}---including the low-dimensional choice-potent subspace embedded within it---is conserved. CKA tests whether two regions share the same \emph{task-specific coordinates} (same directions of variance); Procrustes tests whether they share the same \emph{universal manifold structure} (same geometric shape). High Procrustes with low CKA means: same universal manifold, different subspace sampled within it.

We tested this framework directly by decomposing each region's activity into its LDA choice subspace (potent) and complement (null), computing the fraction of total variance in each. Across 73 regions, LDA places 99.8\% of variance in the null space on average, consistent with the choice signal occupying a tiny fraction of the total variance.

Critically, the geometric dissociation index $\alpha$ is strongly anti-correlated with the null-space fraction: $\rho = -0.68$ ($p = 3 \times 10^{-11}$). Regions where linear and nonlinear metrics disagree most are precisely the regions with the smallest null-space fractions---i.e., the regions where the most variance is concentrated in the choice-potent subspace. Conversely, potent-space decoding accuracy does not correlate with $\alpha$ ($\rho = -0.03$, $p = 0.77$), confirming that the \emph{informativeness} of the potent subspace is constant across regions while its \emph{relative size} drives the metric dissociation. This is exactly the prediction of the potent/null framework: high-dimensional regions have large null spaces that inflate CKA dissimilarity while preserving manifold-level structure.

\subsection{Dimensionality as Computational Strategy}

The task-difficulty modulation ($\Delta\alpha = +12.4$, $p = 3 \times 10^{-7}$) while preserving manifold shape demonstrates that dimensionality is an active computational variable, not a fixed property. Under difficult conditions, circuits compress into fewer dimensions---concentrating variance along task-relevant axes. The evidence-choice subspace analysis reveals two strategies: low-dimensional regions route through shared subspaces (efficient, inflexible), while high-dimensional regions maintain separate, rotatable representations (flexible, requiring nonlinear readout).

\subsection{Toward Neural Taskonomy and RSA of Subspaces}

Our pairwise Grassmannian distances between regions' choice subspaces define a \emph{Grassmannian dissimilarity matrix} (GDM)---an analogue of the representational dissimilarity matrix (RDM) used in RSA \citep{kriegeskorte2008representational}, but operating on subspaces rather than activity patterns. Where RSA asks ``do two conditions evoke similar population responses?'', the GDM asks ``do two regions use similar geometric strategies to encode choice?'' Correlating the GDM with CKA-based and Procrustes-based RDMs, anatomical distance, and functional connectivity would reveal whether subspace similarity captures biologically meaningful organization that standard metrics miss.

This extends naturally to a \emph{neural taskonomy} \citep{zamir2018taskonomy}. Just as Zamir et al.\ discovered that visual tasks form a computational taxonomy with asymmetric transfer, the GDM may reveal a functional hierarchy in which sensory regions' subspaces are close to each other but distant from motor regions', with association regions in between. The structure of this graph---and whether it predicts which region's lesion most disrupts another region's decoding---is a direct test of the manifold framework applied to brain-wide circuit organization.

We constructed such a GDM by computing per-session Grassmannian distances between all co-recorded region pairs (projecting into a shared 30-dimensional PCA ambient space per session), then averaging across sessions. Coverage spans 719 of 2,628 possible region pairs (mean 1.9 sessions per pair).

Two striking patterns emerge. \textbf{First}, the GDM is strongly anti-correlated with CKA similarity ($\rho = -0.96$). Regions that appear similar in CKA's kernel space have \emph{dissimilar} choice subspaces, reinforcing the dissociation: CKA captures shared variance structure (dominated by null-space dimensions) while the GDM captures shared choice geometry (the potent subspace). \textbf{Second}, the GDM does not recover traditional functional groupings (Appendix~\ref{app:gdm_hierarchy})---within-group Grassmannian distances are not smaller than between-group distances---suggesting that choice subspace geometry is organized along a dimension orthogonal to the classical sensory--association--motor hierarchy.

\subsection{Practical Implications}

For neural data analysis, our findings imply a simple diagnostic: compute the power-law exponent $\alpha$ before selecting analysis methods. For $\alpha > 3$ (low-dimensional), linear methods (CKA, LDA) are informative and metrics will generally agree. For $\alpha < 1$ (high-dimensional), nonlinear methods are necessary---and linear methods may be actively misleading, since the large null space will dominate kernel-based comparisons. For intermediate regimes, both should be reported alongside $\alpha$, so readers can assess the degree to which metric choice influences the conclusion.

\subsection{Limitations}

\textbf{Dataset scale.} The Steinmetz dataset comprises 10 mice and 73 regions. The IBL replication (\S\ref{sec:ibl}) confirms the core anti-correlation on an independent dataset (139 mice, 12 labs, 11 matched regions), but absolute power-law exponents do not transfer ($\rho = -0.27$, n.s.), limiting cross-dataset dimensionality comparisons. Per-region significance testing would require denser coverage than the current 3--5 sessions per region in IBL.

\textbf{No temporal resolution.} All analyses use trial-averaged activity (150--350ms window). The metric dissociation may evolve over the trial epoch, but characterizing this requires time-resolved metric computation with sufficient trial counts per time bin.

\textbf{Task specificity.} All results are for binary visual choice. Other task types (navigation, working memory) may show different dimensionality--metric interactions.

\section{Conclusion}

We establish that the choice of geometric similarity metric is not neutral: linear (CKA) and nonlinear (UMAP Procrustes) similarity are strongly anti-correlated across brain regions ($\rho = -0.85$), and this dissociation is fully mediated by effective dimensionality. Controlling for the power-law exponent $\alpha$ reverses the sign, showing that the metrics agree once the spectral structure is accounted for. Singular value spectra are universal across animals while temporal modes are idiosyncratic, and dimensionality itself is an active computational variable---compressing under hard task conditions while preserving manifold shape. Evidence and choice subspaces reveal two encoding strategies: shared-subspace routing in low-dimensional regions versus separate, rotatable representations in high-dimensional regions. A potent/null space decomposition confirms the mechanism: the geometric dissociation index is strongly anti-correlated with null-space fraction ($\rho = -0.86$), meaning regions where metrics disagree most are those with the largest choice-potent subspaces. The CKA--Procrustes dissociation replicates on the independent IBL Brain-Wide Map dataset ($\rho = -0.94$, 139 mice, 12 labs), confirming that this is a property of neural population geometry rather than an artifact of a single recording pipeline. These findings imply a practical diagnostic---report $\alpha$ alongside any similarity metric---and motivate careful attention to the interaction between metric family and dimensionality in any brain-wide comparison. A companion paper extends these findings to causal subspace discovery, showing that the dimensionality interaction also determines which methods can identify task-relevant subspaces.

\bibliographystyle{tmlr}
\bibliography{references}

\appendix
\section{Sequential Discovery and Multiple Comparisons}
\label{app:sequential}

\begin{table}[h]
\centering
\small
\begin{tabular}{@{}clcccc@{}}
\toprule
\textbf{Stage} & \textbf{Test} & \textbf{$\rho$} & \textbf{$p$} & \textbf{$N_k$} & \textbf{Survives} \\
\midrule
\multicolumn{6}{l}{\emph{Stage 1: Core metric comparison}} \\
1 & CKA--Procrustes anti-correlation & $-0.85$ & $\approx 0$ & 1 & \checkmark \\
\midrule
\multicolumn{6}{l}{\emph{Stage 2: Mechanism identification}} \\
2 & Dimensionality mediates (partial $r$) & $+0.44$ & $< 10^{-6}$ & 2 & \checkmark \\
\midrule
\multicolumn{6}{l}{\emph{Stage 3: Spectral and subspace structure}} \\
3 & Spectral universality (mean $r$) & $+0.94$ & $\approx 0$ & 3 & \checkmark \\
4 & Task-difficulty dimensionality shift & $+12.4$ & $3 \times 10^{-7}$ & 4 & \checkmark \\
5 & Evidence-choice alignment vs.\ $\alpha$ & $-0.59$ & $< 10^{-6}$ & 5 & \checkmark \\
\midrule
\multicolumn{6}{l}{\emph{Stage 4: Replication}} \\
6 & IBL CKA--Procrustes anti-correlation & $-0.94$ & $5.5 \times 10^{-5}$ & 6 & \checkmark \\
\midrule
\multicolumn{6}{l}{\emph{Stage 5: Potent/null mechanism}} \\
7 & Dissociation index vs.\ null-space fraction & $-0.86$ & $5 \times 10^{-22}$ & 7 & \checkmark \\
\bottomrule
\end{tabular}
\caption{Sequential discovery table. All seven core claims survive Holm correction.}
\label{tab:sequential}
\end{table}

\section{Robustness Controls}
\label{app:robustness}

\begin{table}[h]
\centering
\small
\caption{Robustness of the CKA--Procrustes anti-correlation.}
\label{tab:robustness}
\begin{tabular}{@{}lccl@{}}
\toprule
\textbf{Control} & \textbf{$\rho$} & \textbf{$p$} & \textbf{$n$} \\
\midrule
Raw (full data) & $-0.850$ & $\approx 0$ & 1,316 pairs \\
Bootstrap 95\% CI & $[-0.867, -0.830]$ & --- & 10,000 resamples \\
Leave-one-mouse-out (range) & $[-0.872, -0.835]$ & all sig. & 10 mice \\
Leave-one-region-out (range) & $[-0.863, -0.841]$ & all sig. & 50 regions \\
Neuron sub-sampling & $-0.814$ & $< 10^{-7}$ & 29 pairs \\
Spectral denoising (MP) & $-0.889$ & $\approx 0$ & 1,316 pairs \\
Trial-label shuffling & $-0.395$ & $< 10^{-50}$ & 1,316 pairs \\
\midrule
\multicolumn{4}{l}{\emph{Partial correlations (Pearson $r$)}} \\
Controlling neuron count & $+0.134$ & --- & sign reversal \\
Controlling $\alpha$ & $+0.437$ & --- & sign reversal \\
Controlling both & $+0.196$ & --- & sign reversal \\
\bottomrule
\end{tabular}
\end{table}

\section{Metric Family Comparison}
\label{app:metrics}

\begin{table}[h]
\centering
\small
\caption{Correlation of each metric family with dimensionality ($\alpha$) and with CKA/Procrustes.}
\label{tab:metrics}
\begin{tabular}{@{}lccccc@{}}
\toprule
\textbf{Metric} & \textbf{vs.\ $\alpha$} & \textbf{$p$} & \textbf{vs.\ CKA} & \textbf{vs.\ Procrustes} \\
\midrule
CKA (linear) & $-0.68$ & $< 10^{-4}$ & --- & $-0.85$ \\
UMAP Procrustes & $+0.68$ & $< 10^{-4}$ & $-0.85$ & --- \\
Sliced Wasserstein & $+0.02$ & $0.47$ & $+0.10$ & $-0.09$ \\
SPD Riemannian & $+0.44$ & $< 10^{-4}$ & --- & --- \\
Fisher-Rao & $+0.40$ & $< 10^{-4}$ & --- & --- \\
\bottomrule
\end{tabular}
\end{table}

\section{Evidence-Choice Subspace Alignment}
\label{app:alignment}

\begin{table}[h]
\centering
\small
\caption{Top and bottom 5 regions by evidence-choice subspace alignment.}
\label{tab:routing}
\begin{tabular}{@{}lccc@{}}
\toprule
\textbf{Region} & \textbf{$d_G(V_{\text{ev}}, V_{\text{ch}})$} & \textbf{$\alpha$} & \textbf{Interpretation} \\
\midrule
\multicolumn{4}{l}{\emph{Highest alignment (shared subspace)}} \\
OT & 0.197 & 104.2 & evidence-choice routing \\
IC & 0.409 & 13.1 & sensorimotor integration \\
MEA & 0.453 & 3.0 & evidence-choice routing \\
LH & 0.471 & 3.9 & evidence-choice routing \\
SCm & 0.486 & 15.6 & motor output \\
\midrule
\multicolumn{4}{l}{\emph{Lowest alignment (separated subspaces)}} \\
VAL & 1.163 & 1.7 & thalamic relay \\
LSc & 1.216 & 2.4 & basal ganglia \\
COA & 1.313 & 0.7 & amygdala \\
CL & 1.381 & 27.7 & thalamic relay \\
EPd & 1.397 & 43.6 & amygdala/piriform \\
\bottomrule
\end{tabular}
\end{table}

\section{Topological and Dynamical Invariants}
\label{app:topo_dyn}

\textbf{Persistent homology.} Only 9/73 regions have $\beta_1 > 0$. Grassmannian vs.\ topological Wasserstein $H_1$ distance: $\rho = 0.25$, $p = 0.52$ ($n = 9$).

\textbf{Rotational dynamics.} 72/73 regions show rotation above null. Mean strength $= 0.053 \pm 0.060$; range $0.008$--$0.40$.

\textbf{Spectral universality.} Singular value spectra: mean $r = 0.94 \pm 0.06$ (1,316 pairs). Temporal modes: mean cosine $= 0.09 \pm 0.07$.

\section{Grassmannian Distance and Functional Hierarchy}
\label{app:gdm_hierarchy}

The GDM does not recover traditional functional groupings. We assigned 39 of 73 regions to broad functional groups (visual, motor, somatosensory, auditory, thalamic, hippocampal, prefrontal) based on the Allen Brain Atlas, yielding 97 within-group and 644 between-group pairs. Within-group Grassmannian distances are not smaller than between-group distances (LDA: Mann-Whitney $U = 34{,}509$, $p = 0.97$). Spectral clustering of the GDM into $k = 7$ groups yields weak agreement with functional labels (ARI $= 0.07$, NMI $= 0.18$). This null result suggests that choice subspace geometry is organized along axes orthogonal to the classical sensory--association--motor hierarchy---consistent with the view that the choice signal is routed through a distributed network whose geometry does not respect anatomical proximity or functional category.

\section{Tucker Decomposition}
\label{app:tucker}

Tucker decomposition on (trials $\times$ neurons $\times$ time) tensors: time factors are most conserved (mean similarity $0.174$), followed by trial ($0.106$) and neuron ($0.067$). Tucker $R^2$ correlates modestly with $\alpha$ ($\rho = 0.275$).

\end{document}
